%%%%%%%%%%%%%%%%%%%%%%%%%%%%%%%%%
%
%       EXERCÍCIO
%
%%%%%%%%%%%%%%%%%%%%%%%%%%%%%%%%%

\definevalorquestao

\begin{exercicioBanco}[\valorquestao]
A massa restante de uma substância radioativa, em gramas, após um tempo \(t\) (em anos), é modelada pela função \(M(t) = 512 \cdot 2^{-0,25t}\). O tempo \(t\), em anos, necessário para que a massa dessa substância seja de \(32\) gramas é igual a:

\newcommand{\alternativas}{%
    \begin{center}
        \begin{tabularx}{\textwidth}{XXXXX}
            (a) \(8\). &
            (b) \(12\). &
            (c) \(16\). &
            (d) \(20\). &
            (e) \(24\).
        \end{tabularx}
    \end{center}
}

\newcommand{\resposta}{C}

% Lógica condicional para exibição
\ifdefstring{\atividade}{lista}{%
    \alternativas
    \vspace{0.5em}
    
    \noindent\textbf{Resposta:} letra \textbf{\resposta}.
}{%
    \ifdefstring{\modo}{objetivo}{%
        \alternativas
    }{%
        % modo = discursiva → não mostra alternativas
    }
}
\end{exercicioBanco}

